% le caratteristiche richieste dall'università sono elencate qui: https://stem.elearning.unipd.it/mod/book/view.php?id=234&chapterid=46#modalita
% 12pt: font richiesto dall'università
% twoside: i margini interni ed esterni sono scambiati per le pagine "a sinistra" e "a destra"
% openright: i capitoli cominciano in pagine dispari ("a destra")
% extreport: supporta 12pt
\documentclass[12pt,a4paper,twoside,openright]{extreport}

\usepackage{amsmath}                            % per avere più controllo sulle equazioni
\usepackage{csquotes}                           % per le citazioni
\usepackage{enumitem}                           % per avere più controllo sulle enumerazioni
\usepackage[
    a4paper,
    top=2cm,bottom=2cm,
    outer=2cm,inner=3cm,
    includeheadfoot
]{geometry}                                     % margini richiesti dall'università
\usepackage{graphicx}                           % per le immagini
\usepackage{icomma}                             % per separare le cifre decimali con una virgola
\usepackage{listings}                             % per il codice con la colorazione della sintassi
\usepackage[a-1a]{pdfx}                         % formato richiesto dall'università
\usepackage[output-decimal-marker={,}]{siunitx} % per le unità di misura
\usepackage{subcaption}                         % per le sottodidascalie
\usepackage{svg}                                % per inserire file SVG
\usepackage{mathtools}

\usepackage[english]{babel}
\selectlanguage{english}

\usepackage[backend=bibtex,style=ieee]{biblatex}
\bibliography{bibliography}

\usepackage{setspace}
\onehalfspacing % interlinea richiesta dall'università

\sloppy % per evitare che il testo in \verb finisca oltre i margini

% questi valori vengono usati nella composizione del frontespizio
\title{Design and Evaluation of a Novel Heuristic for Delete-Free AI Planning}
\author{Andrea Stocco}
\date{GG/MM/AAAA}
\newcommand{\supervisor}{Prof. Domenico Salvagnin}

\usepackage{amsthm}
\theoremstyle{definition}
\newtheorem{definition}{Definition}[section]

\begin{document}
\pagenumbering{Roman}
\pagestyle{empty} % per le prime pagine, non mostrare il numero di pagina

\begin{titlepage}
	% solo il frontespizio deve essere simmetrico rispetto ai margini interno ed esterno
	\newgeometry{hmargin=2.5cm,vmargin=2cm}
	\begin{figure}
		\centering
		\begin{subfigure}[b]{0.4\textwidth}
			\includegraphics[width=\textwidth]{images/logo_unipd}
		\end{subfigure}
		\hfill
		\begin{subfigure}[b]{0.3\textwidth}
			\includegraphics[width=\textwidth]{images/logo_dei}
		\end{subfigure}
	\end{figure}

	\vspace*{\stretch{0.5}}

	\begin{center}
		\makeatletter % serve per poter usare \@...

		% NOTA: il Times New Roman non supporta il maiuscoletto.
		\textsc{Dipartimento di Ingegneria dell'Informazione}\\
		\vspace*{\stretch{0.1}}
		\textsc{Master Degree in Computer Engineering}

		\vspace*{\stretch{0.5}}
		\LARGE
		\textbf{\@title}

		\vspace*{\stretch{1}}
		\normalsize
		\begin{tabular*}{\textwidth}{l @{\extracolsep{\fill}} r}
			\textbf{Supervisor} & \textbf{Master Candidate} \\
			\supervisor       & \@author           \\
		\end{tabular*}

		\vspace*{\stretch{2}}
		\textsc{Academic Year 2024-2025} \\
		\vspace*{\stretch{0.1}}
		% Graduation Date\@date

		\makeatother % serve dopo \makeatletter
	\end{center}
	\restoregeometry
\end{titlepage}

\cleardoublepage

% NOTA: l'ambiente \abstract rimuove il numero della pagina e resetta il contatore delle pagine.
\begin{abstract}
	In this thesis, we present a new heuristic algorithm for solving delete-free AI planning problems.
Delete-free planning, where actions do not remove facts from the world state, offers a simplified yet
meaningful subset of classical planning. We propose an algorithm that estimates goal distances by
leveraging structural properties of delete-free domains. We compare its performance against
existing heuristic approaches on a range of benchmark problems. The evaluation demonstrates that our
method achieves competitive or improved results in terms of planning time and solution quality,
highlighting its potential for use in practical planning systems.

\end{abstract}
\cleardoublepage

\vspace*{\stretch{1}}
\begin{flushright}
	\textit{To my family.}
\end{flushright}
\vspace{\stretch{4}}

\cleardoublepage

\pagestyle{plain} % comincia a mostrare il numero di pagina

\tableofcontents
\cleardoublepage

\listoffigures
\cleardoublepage % per assicurarsi che la numerazione araba cominci col primo capitolo

\listoftables
\cleardoublepage

\pagenumbering{arabic}

\chapter{Introduction}
\label{ch:intro}
\textit{Automated Planning} (or AI Planning) is a core subfield of Artificial Intelligence.
A planning problem, in general terms, can be defined as the task of finding a sequence
of actions that leads from a given initial state to a desired goal state.
While multiple formal models exist for defining planning problems, this thesis
focuses on the \textit{$SAS^+$ formalism} defined in Definition \ref{def:sas+}

\begin{definition}[$\text{SAS}^+$ Planning Task]
	\label{def:sas+}
	A $\text{SAS}^+$ planning task is a 5-tuple $\Pi = \langle V, s_0, s_*, A \rangle$ with
	the following components:
	\begin{itemize}
		\item \(V\): finite set of state variables \(v\),
		      each with finite domain \textit{dom(v)}
		\item \(s_0\): variable assignment defining the initial state
		\item \(s_*\): partial variable assignment defining the goal
		\item \(A\): finite set of actions (or operators),
		      where each action $a \in A$ has the following components:
		      \begin{itemize}
			      \item Preconditions \textit{pre(a)}: partial variable assignment
			      \item Effects \textit{eff(a)}: partial variable assignment
			      \item Cost \textit{cost(a)}: non-negative real number
		      \end{itemize}
	\end{itemize}
\end{definition}

For an action $a \in A$ to be applicable, all of its preconditions
\textit{pre(a)} must be satisfied in the current state. Once applied, its effects
\textit{eff(a)} are used to update the current state accordingly.
In the simplified class of problems considered in this thesis—referred to as \textit{Delete-Free}—actions do not have any negative effects.
This means that once a fact becomes true in a state, it remains true in all subsequent states.
The solutions to planning problems are called \textit{plans}, and they are formally defined in Definition \ref{def:plan}

\begin{definition}[Plan]
	\label{def:plan}
	A plan for a planning problem is a sequence of actions occurring as labels on a path
	from the initial state to a goal state.
	The cost of a plan $\langle a_1, a_2, \dots, a_n \rangle$ is $\sum_{i = 1}^n cost(a_i)$.
	A plan is optimal if it has minimal cost.
\end{definition}

There exist optimal planners, such as \textit{Fast-Downward}, that are guaranteed to find the optimal solution to a given problem.
In contrast, this thesis focuses on \textit{heuristic} algorithms which—by design— do not
guarantee optimality, but aim to efficiently find near-optimal solutions in significantly less time.

TODO

\cleardoublepage

\chapter{Methodology}
\label{ch:methodology}
This chapter outlines the methodological approach used to compare the selected algorithms.
Several algorithms and variations were implemented in this thesis. Some
were chosen solely as baselines for the experiments, such as Random and Greedy.
As their names suggest, the former selects a random action at each iteration, while the latter
always chooses the action with the minimum cost.
The $h^{\max}$ heuristic \cite{bonet2001planning} is one of the most important
heuristics in Delete Free AI Planning problems. A revised version of this heuristic, incorporating
a pruning logic, is proposed in this work.
Additionally, a novel heuristic was developed, which assigns to each action the cost of the shortest path to reach the goal
state and then selects the action that is closest to the goal.\\
The experiments were conducted on a testbed containing 3,104 domain-independent instances.
Each algorithm was run on every instance with 10 different random seeds, resulting in 31,040 runs
per algorithm. A time limit of 60 seconds was set for each run.
All experiments presented in this thesis were conducted on a computing cluster maintained by the Operations Research group at the
Department of Information Engineering (DEI), University of Padova.
The cluster consists of 15 identical blades, each equipped with an Intel(R) Xeon(R) CPU E5-2623 v3 @ 3.00GHz, 16GB of RAM, and running Fedora 37.
The algorithms were compared in terms of solution cost, where a lower cost indicates a better solution.
To evaluate the effectiveness of the algorithms, a \textit{cumulative distribution plot} is used.
The concept of the primal gap, originally defined in the context of Mixed Integer Programming (MIP)
(see Definition \ref{def:primalgap}) \cite{berthold2013measuring}, is adapted here to quantify the quality of solutions across algorithms.

\begin{definition}[Primal Gap]
	\label{def:primalgap}
	Let \~x be a solution for a problem, and $\text{\~x}_{opt}$ be an optimal (or best known)
	solution for that problem. We define the primal gap $\gamma \in \left[0, 1\right]$ of \~x as follows:
	\begin{equation*}
		\gamma\left(\text{\~x}\right)\coloneqq\begin{cases}
			0,                                                                                                                            & \text{if $\lvert c^T\tilde{x}_{opt}\rvert$ = $\lvert c^T\tilde{x}\rvert$ = 0}, \\
			1,                                                                                                                            & \text{if $c^T\tilde{x}_{opt} \cdot c^T\tilde{x} < 0$},                         \\
			\frac{\lvert c^T\tilde{x}_{opt} - c^T\tilde{x} \rvert}{max \{\lvert c^T\tilde{x}_{opt} \rvert, \lvert c^T\tilde{x} \rvert\}}, & \text{else}.
		\end{cases}
	\end{equation*}
\end{definition}

An algorithm that finds the best solution for a given instance will exhibit a primal gap of zero.
For all other algorithms, the primal gap approaches zero as the solution quality approaches the best known,
and approaches one as it deviates further from it.
For each algorithm, a cumulative distribution plot is generated by plotting the percentage of instances
on which the algorithm has a primal gap below a given threshold. A more detailed explanation of this analysis will be provided in the following sections.

\cleardoublepage

\nocite{planning-tutorial}
\printbibliography[heading=bibintoc]
\end{document}
